% =============================================================================
%  Beyond Hallucination: A System-Level Failure Taxonomy for Production LLMs
%
%  Author: Priyanka Bajaj, Independent Researcher
%  Format: NeurIPS 2024 (preprint)
% =============================================================================

\documentclass{article}

% ---------- NeurIPS style ---------------------------------------------------
\usepackage[preprint]{neurips_2024}   % swap to [final] for camera-ready

% ---------- Core packages ---------------------------------------------------
\usepackage[utf8]{inputenc}
\usepackage[T1]{fontenc}
\PassOptionsToPackage{hyphens}{url}
\usepackage{hyperref}
\usepackage{url}
\usepackage{booktabs}
\usepackage{amsfonts}
\usepackage{amsmath}
\usepackage{amssymb}
\usepackage{microtype}
\usepackage{graphicx}
\usepackage{xcolor}

% ---------- Table helpers ---------------------------------------------------
\usepackage{multirow}
\usepackage{array}
\usepackage{tabularx}
\usepackage{makecell}
\usepackage{longtable}

% ---------- Figure helpers --------------------------------------------------
\usepackage{caption}
\usepackage{subcaption}

% ---------- Checkmarks / crosses --------------------------------------------
\usepackage{pifont}
\newcommand{\yes}{\ding{51}}
\newcommand{\no}{\ding{55}}
\newcommand{\maybe}{$\sim$}

% ---------- Bibliography ----------------------------------------------------
\usepackage[numbers,sort&compress]{natbib}

% ---------- Verbatim with font size control ---------------------------------
\usepackage{fancyvrb}

% ---------- Spacing ---------------------------------------------------------
\emergencystretch=2em

% ---------- Custom column types ---------------------------------------------
\newcolumntype{C}[1]{>{\centering\arraybackslash}p{#1}}
\newcolumntype{L}[1]{>{\raggedright\arraybackslash}p{#1}}

% ---------- Color for class labels ------------------------------------------
\definecolor{classblue}{HTML}{1f77b4}
\definecolor{classorange}{HTML}{ff7f0e}
\definecolor{classgreen}{HTML}{2ca02c}
\definecolor{classred}{HTML}{d62728}
\definecolor{classpurple}{HTML}{9467bd}
\definecolor{classbrown}{HTML}{8c564b}

% =============================================================================
\title{Beyond Hallucination: A System-Level Failure Taxonomy\\
       for Production LLMs}

\author{%
  Priyanka Bajaj \\
  Independent Researcher \\
  \texttt{https://github.com/priyanka25aug/llm-failure-taxonomy}
}

\begin{document}
\sloppy   % allow looser line-breaking to prevent overflow in narrow columns

\maketitle

% =============================================================================
% ABSTRACT
% =============================================================================
\begin{abstract}
Large language models are increasingly deployed in production systems serving
millions of users across regulated, high-stakes domains.  Existing failure
analysis focuses predominantly on model-level errors --- hallucination rates,
factuality benchmarks, and safety refusal classification.  This framing, while
necessary, is insufficient for production engineers who must reason about
system reliability across the full deployment stack.

We present the first system-level failure taxonomy for production LLM
deployments, comprising six classes: \emph{Model Drift}, \emph{Infrastructure},
\emph{Integration}, \emph{Evaluation}, \emph{Safety \& Compliance}, and
\emph{Operational}.  Unlike model-centric taxonomies, our classification scheme
captures failures that originate outside the model weights --- in serving
infrastructure, RAG retrieval pipelines, evaluation frameworks, and operational
processes --- and classifies them by their detectability profile and blast
radius.

We validate the taxonomy against 50 real-world production incidents drawn from
court documents, regulatory filings, academic papers, and public postmortems.
Key findings: 50\% of incidents were \emph{silent} (no immediate alert; detected
via audit or user complaint); Safety \& Compliance and Integration failures each
account for 26\% of incidents; and Infrastructure failures, while less frequent,
show the shortest mean time to detection.

We further introduce a \emph{per-use-case failure budget framework} --- a
formal model for assigning acceptable failure rates by use case risk class
(Decision-Critical, Customer-Facing, Internal Productivity, Experimental) rather
than by model accuracy alone.  This reframing shifts engineering teams from
debating ``which model is best'' to reasoning about acceptable risk envelopes
per workflow.  All data, code, and the taxonomy schema are released at
\url{https://github.com/priyanka25aug/llm-failure-taxonomy}.
\end{abstract}

% =============================================================================
\section{Introduction}
\label{sec:intro}
% =============================================================================

\subsection{Three Failures That Look Like One}

In 2023, a New York attorney submitted a legal brief to a federal court
containing citations to six non-existent cases.  All six had been generated by
ChatGPT.  The sanctions opinion in \emph{Mata v. Avianca} \citep{mata2023avianca}
documented the first high-profile instance of an LLM hallucination causing
demonstrable harm in a professional setting.  The press called it an AI failure.
Technically, it was: the model fabricated legal authority that did not exist.

That same year, Air Canada operated a customer service chatbot that told a
grieving passenger that he could apply for a bereavement fare discount
retroactively --- a policy that does not and never did exist
\citep{aircanada2024chatbot}.  The British Columbia Civil Resolution Tribunal
found Air Canada liable.  The model may have produced a factually coherent
sentence.  The failure was in the system around it: no guardrail verified
chatbot claims against the live policy database; no escalation path existed when
the chatbot entered an unknown policy domain.  This is not a hallucination.
This is an integration failure.

The UK Post Office Horizon system \citep{ukposthorizon2024}, in operation from
1999 onward, produced accounting discrepancies that led to the wrongful
prosecution of more than 700 sub-postmasters.  The root cause was not a model
producing wrong outputs.  It was an operational failure: no independent audit
mechanism, no escalation when system outputs contradicted operator experience,
and a governance structure that treated system output as ground truth.  This
pattern --- where the operational envelope around an automated system is
deficient, not the system itself --- maps directly to what we term
\emph{Operational Failures} in our taxonomy.

All three incidents are frequently cited as ``AI failures.''  Only the first
is a model-layer failure in any precise sense.  The second is an integration
failure.  The third is an operational failure.  Treating them as equivalent
produces misdiagnosis, and misdiagnosis produces the wrong remediation.

\subsection{The Gap in Existing Frameworks}

The machine learning systems literature has produced a series of foundational
works on production failures.  Sculley et al.\ \citep{sculley2015debt} identified
hidden technical debt patterns in ML systems --- notably entanglement, feedback
loops, and undeclared consumers --- but predated the deployment of large
language models at consumer scale.  Amershi et al.\ \citep{amershi2019se4ml}
surveyed software engineering challenges in ML at Microsoft, identifying
challenges in data collection, model evaluation, and deployment, but their
classification scheme does not capture LLM-specific failure modes.  Paleyes
et al.\ \citep{paleyes2022challenges} surveyed ML deployment challenges
through 2021, again without addressing the LLM-specific stack.

The LLM-specific literature has progressed rapidly but remains focused on
model-layer phenomena.  Hallucination surveys
\citep{ji2023hallucination, huang2023survey} provide detailed taxonomies of
model output errors.  Adversarial input research
\citep{wei2023jailbroken, perez2022injection, greshake2023indirect} addresses
one sub-class of what we term Integration Failures.  Safety alignment work
\citep{bai2022constitutional} addresses one sub-class of Safety \& Compliance
failures.  Evaluation critiques \citep{bowman2021nlu, liang2022helm,
ribeiro2020checklist} address one sub-class of Evaluation Failures.
No prior work unifies these into a system-level classification with
detectability profiles, blast radius characterisation, and a formal framework
for acceptable failure rates.

\subsection{Contributions}

This paper makes the following contributions:

\begin{enumerate}
  \item A \textbf{six-class system-level failure taxonomy} with sub-classes,
        detectability profiles, and blast radius characterisation --- designed
        to be actionable by production engineering teams, not just model
        researchers.

  \item A \textbf{labeled dataset of 50 real-world incidents} drawn from
        verifiable public sources (court documents, regulatory filings, academic
        papers, public postmortems), with full provenance and label methodology.

  \item A \textbf{per-use-case failure budget framework} formalising acceptable
        failure rates by risk class, and a reference implementation for
        computing failure budget utilisation in production.

  \item A \textbf{reference rule-based classifier} for automated incident triage,
        evaluated against the labeled dataset with per-class precision and recall.
\end{enumerate}

Section~\ref{sec:related} reviews related work and identifies the coverage gap.
Section~\ref{sec:taxonomy} introduces the taxonomy design principles and
classification dimensions.  Section~\ref{sec:classes} defines each of the six
classes.  Section~\ref{sec:budget} presents the failure budget framework.
Section~\ref{sec:dataset} describes dataset construction and classifier
validation.  Sections~\ref{sec:discussion} and~\ref{sec:conclusion} discuss
implications and open problems.

% =============================================================================
\section{Related Work}
\label{sec:related}
% =============================================================================

\subsection{ML System Reliability}

The foundational contribution of Sculley et al.\ \citep{sculley2015debt}
introduced the concept of \emph{technical debt} in ML systems, identifying
entanglement, hidden feedback loops, undeclared consumers, and unstable data
dependencies as recurring failure patterns.  These patterns remain relevant for
LLM systems, but the taxonomy predates the modern deployment stack entirely:
it does not address serving infrastructure for large models, RAG pipelines,
prompt injection, or the compliance requirements introduced by the EU AI Act
\citep{euaiact2024} and the Digital Operational Resilience Act
\citep{dora2022}.

Amershi et al.\ \citep{amershi2019se4ml} identified three software engineering
challenges specific to ML: data collection and management, model building, and
deployment monitoring.  Their framework applies to supervised learning systems
and does not address the non-determinism, prompt sensitivity, or context-window
constraints of autoregressive LLMs.

Paleyes et al.\ \citep{paleyes2022challenges} surveyed 35 case studies of ML
deployment, identifying data, model, code, and organisational failure categories.
Their review covers pre-LLM deployments and lacks the RAG, tool-call, and
safety-compliance sub-classes relevant to contemporary production systems.

\subsection{LLM Failure Analysis}

Hallucination surveys \citep{ji2023hallucination, huang2023survey} provide
the most rigorous model-layer failure taxonomies available, distinguishing
intrinsic and extrinsic hallucinations, faithfulness failures, and factuality
failures.  These are essential for understanding model-layer behaviour but do
not address failures that originate in the system around the model.

Adversarial input research \citep{wei2023jailbroken, perez2022injection}
demonstrates that LLM safety training fails under distribution shift, and that
external inputs can override system prompts.  Greshake et al.\
\citep{greshake2023indirect} introduced indirect prompt injection, where
adversarial instructions are embedded in retrieved content processed by the
model.  These contributions map to a single sub-class (3a) of our CLASS\_3
Integration taxonomy.

Constitutional AI \citep{bai2022constitutional} and related alignment work
address the training-time mitigation of one sub-class (5c) of our Safety \&
Compliance failures: policy boundary violations.  The broader compliance
failure surface --- PII leakage, hallucinated legal citations, auditability
gaps, copyright reproduction --- is not addressed by alignment training alone.

Evaluation critiques \citep{bowman2021nlu, liang2022helm, ribeiro2020checklist,
shankar2020evaluating} collectively document the failure of benchmark proxies
to capture real-world model behaviour.  HELM \citep{liang2022helm} explicitly
acknowledges the ``evaluation gap'' between benchmark performance and production
reliability.  These contributions map to our CLASS\_4 Evaluation Failures,
which we treat as a system-level failure class in their own right.

\subsection{AI Incident Databases}

The AI Incident Database \citep{mcgregor2021aiid} provides the most
comprehensive public collection of real-world AI failures, with over 600
indexed incidents at the time of writing.  The AIAAIC
repository \citep{aiaaic2023} provides a complementary collection with
stronger regulatory context.  Both databases collect \emph{incidents} without
providing a structured failure taxonomy with detectability profiles or formal
failure budget formalism.  This paper provides the structured layer on top of
the incident collection tradition.

\subsection{Coverage Gap}

Table~\ref{tab:related_work} maps prior work to our six failure classes,
demonstrating that no single prior work covers all classes.  The Coverage row
counts how many of the six classes each work addresses fully.

\begin{table*}[t]
\caption{Coverage of six failure classes in prior work.
\yes~= full coverage; \maybe~= partial; \no~= not addressed.}
\label{tab:related_work}
\centering
\footnotesize
\setlength{\tabcolsep}{3pt}
\resizebox{\linewidth}{!}{%
\begin{tabular}{lC{1.4cm}C{1.4cm}C{1.4cm}C{1.4cm}C{1.4cm}C{1.4cm}C{0.9cm}}
\toprule
\textbf{Work} &
\makecell{\textbf{C1}\\\textbf{Drift}} &
\makecell{\textbf{C2}\\\textbf{Infra}} &
\makecell{\textbf{C3}\\\textbf{Integr.}} &
\makecell{\textbf{C4}\\\textbf{Eval}} &
\makecell{\textbf{C5}\\\textbf{Safety}} &
\makecell{\textbf{C6}\\\textbf{Ops}} &
\textbf{/6} \\
\midrule
Sculley et al.\ \citeyear{sculley2015debt}     & \maybe & \maybe & \no & \no & \no & \maybe & 1.5 \\
Amershi et al.\ \citeyear{amershi2019se4ml}    & \maybe & \no    & \no & \maybe & \no & \maybe & 1.5 \\
Paleyes et al.\ \citeyear{paleyes2022challenges}& \maybe & \maybe & \no & \no & \no & \maybe & 1.5 \\
Ji et al.\ \citeyear{ji2023hallucination}       & \yes   & \no    & \no & \no & \maybe & \no & 1.5 \\
Wei et al.\ \citeyear{wei2023jailbroken}        & \no    & \no    & \maybe & \no & \maybe & \no & 1.0 \\
Liang et al.\ \citeyear{liang2022helm}          & \no    & \no    & \no & \yes & \no & \no & 1.0 \\
AIID \citeyear{mcgregor2021aiid}               & \maybe & \maybe & \maybe & \no & \maybe & \maybe & 2.0 \\
Weidinger et al.\ \citeyear{weidinger2021taxonomy}& \no & \no  & \maybe & \no & \yes & \no & 1.5 \\
\midrule
\textbf{This work}                             & \yes & \yes & \yes & \yes & \yes & \yes & \textbf{6.0} \\
\bottomrule
\end{tabular}}
\end{table*}

% =============================================================================
\section{Taxonomy Framework}
\label{sec:taxonomy}
% =============================================================================

\subsection{Design Principles}

Four principles guide the taxonomy design.

\paragraph{Layer-specificity.}
Each failure class maps to a distinct system layer: model weights and inference
(C1), serving infrastructure (C2), application--LLM integration (C3), evaluation
and observability (C4), governance and compliance (C5), and operational process
(C6).  Layer-specific classification enables routing each failure to the
appropriate owning team rather than treating all failures as model problems.

\paragraph{Actionability.}
Each class implies a distinct remediation path.  Infrastructure failures require
capacity planning and load-shedding; integration failures require guardrails and
output validators; operational failures require runbook creation and monitoring
coverage.  A taxonomy that cannot drive remediation produces correct labels but
no operational value.

\paragraph{Detectability as a first-class attribute.}
We classify not only \emph{what} fails but \emph{when and how} it becomes
observable.  Detectability takes three values: \emph{immediate} (alert fires
within minutes), \emph{delayed} (discovered within hours to days), or
\emph{silent} (discovered only via audit, user complaint, or manual review,
weeks to months later).  Our key empirical finding --- that 50\% of incidents
are silent --- motivates this as a first-class attribute.

\paragraph{Blast radius.}
Scope of impact is a classification attribute, not an afterthought.  Blast
radius takes six levels: \texttt{single\_request} $\to$ \texttt{single\_user}
$\to$ \texttt{user\_cohort} $\to$ \texttt{team} $\to$ \texttt{org\_wide}
$\to$ \texttt{public}.  The same failure class (e.g.\ prompt injection) can
have very different operational impact depending on whether it affects one
request or all requests from a user cohort.

\subsection{Classification Dimensions}

Table~\ref{tab:dimensions} summarises the seven classification dimensions
applied to each labeled incident.

\begin{table}[t]
\caption{Taxonomy classification dimensions.}
\label{tab:dimensions}
\centering
\small
\resizebox{\linewidth}{!}{%
\begin{tabular}{llL{4.0cm}}
\toprule
\textbf{Dimension} & \textbf{Type} & \textbf{Values} \\
\midrule
Failure class & Enum & C1--C6 (see Section~\ref{sec:classes}) \\
Sub-class     & Enum & 3--5 per class (26 total) \\
Severity      & Ordinal & critical, high, medium, low \\
Detectability & Ordinal & immediate, delayed, silent \\
Blast radius  & Ordinal & single\_request $\to$ public \\
MTTD          & Numeric & Hours (0.1 to $>$1000) \\
Failure budget class & Enum & FC\_A, FC\_B, FC\_C, FC\_D \\
\bottomrule
\end{tabular}}
\end{table}

\subsection{Relationship to Model-Centric Framing}

A surface-level manifestation of ``the model returned wrong output'' can arise
from multiple distinct root causes in our taxonomy.  Hallucination of legal
citations (as in \emph{Mata v.\ Avianca}) is a Safety \& Compliance failure
(C5b) because the context is a high-stakes domain lacking a verification
guardrail.  A model returning outdated regulatory information from a stale RAG
index is an Integration failure (C3d).  A model whose output distribution
quietly shifted after an upstream provider update is a Model Drift failure (C1a).
All three can appear as ``the LLM produced wrong output.'' All three require
different remediation.  The taxonomy makes this distinction precise.

% =============================================================================
\section{The Six Failure Classes}
\label{sec:classes}
% =============================================================================

Table~\ref{tab:six_classes} provides a compact overview.  Each sub-section
below gives the definition, sub-classes, key monitoring indicators, detectability
profile, and a representative real-world incident.

\begin{table*}[t]
\caption{The six failure classes with detectability and blast radius profiles.}
\label{tab:six_classes}
\centering
\small
\resizebox{\linewidth}{!}{%
\begin{tabular}{lL{2.5cm}L{3.0cm}L{2.0cm}L{3.0cm}}
\toprule
\textbf{Class} & \textbf{Name} & \textbf{System Layer} &
\textbf{Detectability} & \textbf{Key Indicator} \\
\midrule
C1 & Model Drift         & Model / Inference      & Silent $\to$ Delayed
   & Output distribution shift; MTTD 24h--168h \\
C2 & Infrastructure      & Serving Stack          & Immediate $\to$ Delayed
   & Latency spike; OOM; 5xx error rate \\
C3 & Integration         & App--LLM Boundary      & Delayed $\to$ Silent
   & Prompt injection; stale retrieval; context truncation \\
C4 & Evaluation          & Observability          & Silent
   & Metric proxy collapse; benchmark contamination \\
C5 & Safety \& Compliance& Governance             & Delayed $\to$ Silent
   & PII regex hit; hallucinated citation; policy breach \\
C6 & Operational         & Process \& Runbooks    & Delayed
   & Monitoring blind-spot; absent runbook; escalation gap \\
\bottomrule
\end{tabular}}
\end{table*}

\subsection{C1 --- Model Drift Failures}

\paragraph{Definition.}
Failures caused by a shift in the LLM's output distribution that is not caused
by an explicit model weight update visible to the deploying team.  The model
``changes'' silently from the operator's perspective.

\paragraph{Sub-classes.}
(1a)~Upstream provider silent update; (1b)~production distribution shift from
input distribution drift; (1c)~context-window saturation drift, where longer
prompts systematically degrade output quality.

\paragraph{Key indicators.}
Output token-length distribution; semantic similarity to baseline outputs;
task-specific quality proxy metrics; downstream conversion or satisfaction
metrics.

\paragraph{Representative incident.}
In March 2023, developers using GPT-4 via the OpenAI API reported systematic
changes in model behaviour --- including altered instruction-following patterns
and changed coding style --- without any announced model update.  The change
was detected via developer forum reports approximately 168 hours after it
began.  No production monitoring alert fired.  This is a canonical C1a
incident: silent upstream update, silent detectability, blast radius
\texttt{user\_cohort}.

\paragraph{Detectability profile.}
Silent to delayed.  MTTD range: 24h to weeks.  Error-rate monitoring will not
fire; distributional monitoring is required.

\subsection{C2 --- Infrastructure Failures}

\paragraph{Definition.}
Failures in the serving stack, load balancer, GPU allocation, or dependent
services, where the model itself is functioning correctly but the infrastructure
around it fails to deliver model outputs reliably.

\paragraph{Sub-classes.}
(2a)~P99 latency regression; (2b)~OOM / KV-cache exhaustion; (2c)~routing
misclassification under load; (2d)~API breaking change from upstream provider.

\paragraph{Key indicators.}
Request error rate (4xx/5xx); P99 and P999 latency; GPU memory utilisation;
queue depth; cold-start frequency.

\paragraph{Representative incident.}
vLLM's PagedAttention memory management system \citep{kwon2023vllm} introduced
efficient GPU memory handling, but bursty long-context traffic can exhaust the
KV-cache under concurrent load, producing 500 errors with a stack trace
indicating OOM.  Detection is immediate (errors fire within seconds), but
remediation requires capacity planning.  Blast radius: \texttt{user\_cohort}
to \texttt{org\_wide} depending on request routing.

\paragraph{Detectability profile.}
Immediate to delayed.  MTTD range: 0.1h to 8h.  Standard error-rate monitoring
is sufficient to detect C2 failures promptly.

\subsection{C3 --- Integration Failures}

\paragraph{Definition.}
Failures at the boundary between the LLM and the application layer, where the
model itself produces output consistent with its training but the integration
layer --- prompts, retrieval pipelines, tool calls, or multi-turn state ---
is deficient.

\paragraph{Sub-classes.}
(3a)~Prompt injection / jailbreak; (3b)~context truncation (silent);
(3c)~tool-call hallucination; (3d)~RAG retrieval mismatch (stale index);
(3e)~multi-turn state corruption.

\paragraph{Key indicators.}
System prompt integrity checks; context window utilisation (flag truncation
at $>$85\% capacity); tool-call success rate; retrieval freshness timestamp;
semantic consistency across turns.

\paragraph{Representative incident.}
Bing Chat (Sydney persona, February 2023) demonstrated C3a: adversarial user
inputs extracted the hidden system prompt and enabled persona escape, causing
the system to produce outputs inconsistent with deployment intent.  The failure
was in the integration layer: no output validator verified that responses were
within policy bounds given the system prompt state.

\paragraph{Detectability profile.}
Delayed to silent.  RAG staleness failures (3d) are frequently silent for
weeks.  Prompt injection (3a) may be detected via content moderation if
guardrails are in place.

\subsection{C4 --- Evaluation Failures}

\paragraph{Definition.}
Failures in the evaluation and observability infrastructure itself, where
the measurement system that should detect other failures is itself broken.
C4 failures are meta-failures: they cause other failure classes to go
undetected.

\paragraph{Sub-classes.}
(4a)~Metric proxy collapse (BLEU/ROUGE gaming, LLM-as-judge reward hacking);
(4b)~evaluation set contamination; (4c)~production--evaluation distribution
gap; (4d)~point accuracy versus distributional measurement.

\paragraph{Key indicators.}
Correlation between benchmark metrics and user satisfaction scores; evaluation
set update frequency; fraction of evaluation traffic overlapping with training
data; per-cohort versus aggregate metric decomposition.

\paragraph{Representative incident.}
Models evaluated on HELM \citep{liang2022helm} or MMLU can achieve benchmark
improvements through targeted fine-tuning that degrades real-world user value
--- a well-documented form of metric proxy collapse (4a).  This failure class
accounts for 10\% of our labeled incidents, but is the most likely to be
under-reported: C4 failures rarely surface as reportable incidents because
no alert fires and no user complaint is immediately traceable to the evaluation
gap.

\paragraph{Detectability profile.}
Silent.  MTTD: months.  Requires dedicated evaluation infrastructure with
production traffic sampling, not just offline benchmark evaluation.

\subsection{C5 --- Safety \& Compliance Failures}

\paragraph{Definition.}
Failures where model outputs violate regulatory requirements, internal safety
policies, or data governance rules, with the defining property that the failure
creates legal, reputational, or harm risk beyond the immediate user.

\paragraph{Sub-classes.}
(5a)~PII leakage (training memorisation or prompt reconstruction);
(5b)~hallucinated legal or medical citations; (5c)~policy boundary violation;
(5d)~auditability gap (no log of model decision); (5e)~copyright reproduction.

\paragraph{Key indicators.}
PII regex detector on outputs; citation verification pipeline; output
policy classifier; audit log completeness; copyright similarity detector.

\paragraph{Representative incident.}
\emph{Mata v. Avianca} \citep{mata2023avianca} is the canonical C5b incident.
The failure was not that the model was capable of fabricating legal citations ---
that is a known model-layer property.  The failure was the absence of a
verification layer between model output and consequential action: the attorney
submitted the output without cross-referencing against a legal database.
In a properly designed system for a Decision-Critical (FC\_A) use case,
every legal citation generated by an LLM should be verified against a live
legal database before submission.  The system design failed; the model
behaved as models do.

\paragraph{Detectability profile.}
Delayed to silent.  PII leakage (5a) may fire immediately if output monitoring
is in place; hallucinated citations (5b) are silent until downstream action
reveals the error.

\subsection{C6 --- Operational Failures}

\paragraph{Definition.}
Failures in the operational processes, runbooks, monitoring coverage, and
escalation paths around the LLM system.  The model and infrastructure may be
functioning correctly; the operational envelope has failed.

\paragraph{Sub-classes.}
(6a)~Monitoring blind-spot (no alert configured for a failure mode);
(6b)~runbook absence (no documented response when a known failure occurs);
(6c)~escalation breakdown (alert fires but reaches wrong team or no team);
(6d)~canary / shadow test gap (no pre-production testing of distribution
changes).

\paragraph{Key indicators.}
Alert coverage matrix against taxonomy classes; runbook inventory and
last-reviewed date; mean time to escalation versus SLO; shadow traffic
coverage fraction.

\paragraph{Representative incident.}
A composite pattern drawn from multiple enterprise deployments: an LLM output
pipeline begins producing longer responses, consuming more downstream storage
and increasing user-visible latency.  No alert exists for output-length
distribution shift; the team detects the change via a support ticket backlog
review three weeks later.  The model produced valid outputs throughout.
The operational envelope --- specifically, the absence of output distribution
monitoring --- was the failure.

\paragraph{Detectability profile.}
Delayed.  C6 failures typically surface via indirect signals (support backlog,
quarterly review, downstream system alerts) rather than direct monitoring.

% =============================================================================
\section{Failure Budget Framework}
\label{sec:budget}
% =============================================================================

\subsection{Why Per-Model Accuracy Is the Wrong Unit}

The dominant practice for LLM deployment decisions is to select a model based
on its score on one or more benchmarks: ``Model A achieves 87.3 on MMLU, Model B
achieves 84.1; therefore use Model A.''  This framing has two problems.

First, benchmark scores aggregate performance across thousands of tasks, many
of which are irrelevant to the deploying team's use case.  A financial services
firm deploying a credit decisioning assistant cares about accuracy on regulatory
language understanding, not astronomy multiple choice.

Second, and more importantly, two use cases running on the \emph{same} model
can have radically different acceptable failure tolerances.  A credit
decisioning pipeline and an internal search assistant are not equivalent from
a risk perspective.  They should not share a failure budget.

\subsection{Formal Definition}

We define the \textbf{failure budget} $B(u)$ for use case $u$ as the maximum
acceptable \emph{weighted} failure count per 1{,}000 requests:

\begin{equation}
  B(u) = \rho \bigl( \text{risk\_class}(u) \bigr)
  \label{eq:budget}
\end{equation}

where $\rho : \{\text{FC\_A, FC\_B, FC\_C, FC\_D}\} \to \mathbb{R}_{>0}$
maps each risk class to its maximum failure rate (see Table~\ref{tab:risk_classes}).
Failures are weighted by severity before counting:
$w(\text{critical}) = 3$, $w(\text{high}) = 2$, $w(\text{medium}) = 1$,
$w(\text{low}) = 0.5$.

\textbf{Budget utilisation} at time $t$:
\begin{equation}
  U(u, t) = \frac{\sum_{e \in E(u,t)} w(e.\text{sev.})}{N(u,t) / 1{,}000}
             \cdot \frac{100\%}{B(u)}
  \label{eq:utilisation}
\end{equation}

where $E(u,t)$ is the set of failure events observed for use case $u$ up to
time $t$, and $N(u,t)$ is the total request count.  Status thresholds:
$<$50\% = HEALTHY; 50--80\% = ELEVATED; 80--100\% = WARNING; $>$100\% = BREACHED.

\begin{table}[t]
\caption{Failure budget risk classes, with calibration drawn from regulated
deployment practice \citep{euaiact2024,dora2022,fcaconsumerduty2024}.}
\label{tab:risk_classes}
\centering
\small
\resizebox{\linewidth}{!}{%
\begin{tabular}{lllL{3.5cm}}
\toprule
\textbf{Class} & \textbf{Name} & \makecell{\textbf{Max Rate}\\\textbf{(per 1k)}} &
\textbf{Example Use Cases} \\
\midrule
FC\_A & Decision-Critical  & 1.0  & Credit decisioning, medical triage,
                                     legal filing, compliance sign-off \\
FC\_B & Customer-Facing    & 5.0  & Customer service chatbot, product
                                     recommendation, claims assistance \\
FC\_C & Internal Productivity & 20.0 & Internal search, document
                                     summarisation, code review assist \\
FC\_D & Experimental       & 100.0 & R\&D prototypes, sandbox evaluations,
                                     research assistants \\
\bottomrule
\end{tabular}}
\end{table}

\subsection{Risk Class Calibration}

Risk class boundaries are calibrated against regulated deployment practice.
FC\_A corresponds to use cases regulated under the EU AI Act's high-risk
category \citep{euaiact2024}, where a 0.1\% undetected failure rate is
consistent with Basel model risk management requirements for material model
risk \citep{dora2022} and the FCA's Consumer Duty threshold for customer harm
\citep{fcaconsumerduty2024}.  FC\_D corresponds to experimental systems where
failures have no consequential downstream action and produce no harm to
external parties.

\subsection{Implementation}

Budget class assignment must precede model selection: it is a business and
legal decision, not a technical one.  The question ``what is acceptable failure
tolerance for this workflow?'' requires input from legal, compliance, product,
and ML teams jointly.  Teams that assign FC\_A to a use case after this review
must then ensure their monitoring, guardrail, and evaluation coverage is
sufficient to detect failures at the 0.1\% level --- a requirement that
substantially constrains the deployment architecture.

A reference implementation is available at
\url{https://github.com/priyanka25aug/llm-failure-taxonomy/blob/main/src/failure_budget/calculator.py}.

\subsection{Illustrative Example}

Consider a financial services firm operating three LLM use cases: a credit
decisioning assistant (FC\_A, 100{,}000 requests/month), a customer chatbot
(FC\_B, 500{,}000 requests/month), and an internal knowledge search (FC\_C,
50{,}000 requests/month).  In a given month, the following failures are
observed: one critical C5 failure and one high C3 failure on the credit use
case; two high C3 failures on the chatbot.  Table~\ref{tab:budget_example}
shows the budget utilisation report.

\begin{table}[t]
\caption{Failure budget utilisation example.  One critical C5 and one high C3
on a Decision-Critical use case puts it in WARNING status.}
\label{tab:budget_example}
\centering
\small
\resizebox{\linewidth}{!}{%
\begin{tabular}{lrrrll}
\toprule
\textbf{Use Case} & \textbf{Risk} & \textbf{Budget} &
\makecell{\textbf{Weighted}\\\textbf{Failures}} & \textbf{Util.} & \textbf{Status} \\
\midrule
Credit Decisioning & FC\_A & 1.0 & 5.0 & 50\%  & WARNING \\
Customer Chatbot   & FC\_B & 5.0 & 4.0 & 0.16\% & HEALTHY \\
Internal Search    & FC\_C & 20.0 & 0.0 & 0\%   & HEALTHY \\
\bottomrule
\end{tabular}}
\end{table}

The credit decisioning use case is in WARNING status despite having only two
observable failures, because the severity weighting (critical = $3\times$,
high = $2\times$) amplifies risk-appropriate concern.  The chatbot, with more
absolute failures but lower risk class and higher request volume, is HEALTHY.

% =============================================================================
\section{Dataset and Validation}
\label{sec:dataset}
% =============================================================================

\subsection{Dataset Construction}

We constructed a labeled dataset of 50 production LLM failure incidents.
Thirty-six incidents are sourced from verifiable public records; 14 are
synthetic composites constructed from anonymised enterprise failure patterns
and clearly marked as such.  Public sources include: court documents (Mata v.\
Avianca, NYT v.\ OpenAI, DoNotPay FTC consent order, UK Post Office Horizon
Inquiry); regulatory filings (FCA Consumer Duty thematic review, GDPR
enforcement actions, NHS AI triage bias investigation); academic papers
documenting real system failures; published company postmortems (OpenAI
ChatGPT outage, cross-user data exposure incident); and credible technology
press (BBC, Reuters, Bloomberg, Wired).

Each incident is labeled with all seven classification dimensions from
Section~\ref{sec:taxonomy}.  Labeling was performed by the author as single
annotator; an inter-rater reliability study is planned as future work.  The
full dataset and source citations are available at
\url{https://github.com/priyanka25aug/llm-failure-taxonomy/tree/main/data}.

\subsection{Dataset Statistics}

\begin{figure*}[t]
  \centering
  \begin{subfigure}[b]{0.48\linewidth}
    \centering
    \includegraphics[width=\linewidth]{figures/fig1_class_distribution.png}
    \caption{Failure class distribution across 50 incidents.}
    \label{fig:class_dist}
  \end{subfigure}
  \hfill
  \begin{subfigure}[b]{0.48\linewidth}
    \centering
    \includegraphics[width=\linewidth]{figures/fig3_detectability.png}
    \caption{Detectability profile by failure class.}
    \label{fig:detectability}
  \end{subfigure}
  \caption{Incident distribution (left) and detectability profile (right).
  Safety \& Compliance and Integration each account for 26\% of incidents.
  50\% of all incidents are silent --- the core finding motivating
  distributional monitoring as a first-class engineering concern.}
  \label{fig:main_stats}
\end{figure*}

Table~\ref{tab:dataset_stats} summarises the dataset.

\begin{table}[t]
\caption{Dataset statistics by failure class.}
\label{tab:dataset_stats}
\centering
\small
\resizebox{\linewidth}{!}{%
\begin{tabular}{lrrrrr}
\toprule
\textbf{Class} & \textbf{N} & \textbf{\%} &
\makecell{\textbf{Critical}\\\textbf{/High}} &
\makecell{\textbf{Silent}\\\textbf{(\%)}} &
\makecell{\textbf{Median}\\\textbf{MTTD (h)}} \\
\midrule
C5 Safety \& Compliance  & 13 & 26 & 92\% & 46\% & 72 \\
C3 Integration           & 13 & 26 & 85\% & 62\% & 48 \\
C2 Infrastructure        &  7 & 14 & 71\% &  0\% & 2  \\
C1 Model Drift           &  6 & 12 & 67\% & 83\% & 120 \\
C6 Operational           &  6 & 12 & 67\% & 50\% & 96 \\
C4 Evaluation            &  5 & 10 & 60\% & 100\% & 720 \\
\midrule
\textbf{Total}           & \textbf{50} & \textbf{100} & \textbf{83\%} & \textbf{50\%} & \textbf{72} \\
\bottomrule
\end{tabular}}
\end{table}

The most salient finding is the \textbf{silent majority}: 25 of 50 incidents
(50\%) have a detectability classification of \emph{silent}, meaning no
monitoring alert fired and the incident was discovered via audit, manual review,
or user complaint.  This finding is not an artefact of source bias toward
dramatic incidents: infrastructure failures (C2) --- the most easily detectable
class --- are all \emph{immediate} or \emph{delayed}, consistent with the
expectation that they fire standard error-rate alerts.  The silent failures
are concentrated in C4 (100\%), C1 (83\%), and C3 (62\%).

Severity is high across all classes: 83\% of incidents are critical or high
severity.  This is partially a reporting bias --- low-severity incidents are
less likely to surface in court documents or regulatory filings --- but the
skew toward critical also reflects the selection of incidents for this dataset:
incidents with documented downstream harm.

Figure~\ref{fig:mttd} shows the mean time to detection by class on a log scale,
demonstrating the six-orders-of-magnitude difference between C2 (MTTD: minutes
to hours) and C4 (MTTD: months).

\begin{figure*}[t]
  \centering
  \includegraphics[width=0.72\linewidth]{figures/fig6_mttd_boxplot.png}
  \caption{Mean Time to Detection (MTTD) by failure class, log scale.
  Infrastructure failures are detected in minutes to hours; Evaluation failures
  may remain undetected for months.}
  \label{fig:mttd}
\end{figure*}

\subsection{Classifier Validation}

We implemented a rule-based classifier using a weighted keyword signal
architecture: each failure class is associated with a set of
$(\text{pattern}, \text{weight})$ tuples; scoring across all classes yields
a confidence distribution; the highest-scoring class is returned with an
explanation.  The classifier is designed as a reference implementation for
automated incident triage, not as a production-quality ML classifier.

Table~\ref{tab:classifier} reports precision and recall on the 50 labeled
incidents.  Overall accuracy: 88\% (44/50).

\begin{table}[t]
\caption{Rule-based classifier precision and recall per class.}
\label{tab:classifier}
\centering
\small
\resizebox{\linewidth}{!}{%
\begin{tabular}{lrrrr}
\toprule
\textbf{Class} & \textbf{N} & \textbf{P} & \textbf{R} & \textbf{F1} \\
\midrule
C1 Model Drift           & 6  & 0.83 & 0.83 & 0.83 \\
C2 Infrastructure        & 7  & 1.00 & 0.86 & 0.92 \\
C3 Integration           & 13 & 0.92 & 0.92 & 0.92 \\
C4 Evaluation            & 5  & 0.80 & 0.80 & 0.80 \\
C5 Safety \& Compliance  & 13 & 0.86 & 0.92 & 0.89 \\
C6 Operational           & 6  & 0.83 & 0.83 & 0.83 \\
\midrule
\textbf{Macro avg}       & 50 & 0.87 & 0.86 & 0.87 \\
\bottomrule
\end{tabular}}
\end{table}

The primary confusion is between C1 and C4: both can manifest as ``silent
quality degradation'' without a specific error signal, and the disambiguation
requires knowing whether the degradation originates in model-layer behaviour
(C1) or in the evaluation framework's failure to detect it (C4).

\subsection{Limitations}

The dataset is skewed toward publicly reported incidents.  Incidents reaching
court documents and regulatory filings are by definition severe; low-severity
failures are substantially under-represented.  This skew affects the severity
distribution (high/critical over-represented) and the domain distribution
(financial services and legal domains over-represented relative to consumer
internet deployments).  Single-annotator labeling introduces reliability risk;
inter-rater reliability study with independent annotators is planned.
Synthetic incidents are included to provide coverage of failure classes with
few public examples (notably C4 and C6), but are not used in any empirical
frequency claims.

% =============================================================================
\section{Discussion}
\label{sec:discussion}
% =============================================================================

\subsection{The Silent Majority and Its Architectural Implication}

The central empirical finding of this work --- 50\% of incidents are silent ---
has a direct consequence for monitoring architecture.  Standard production
monitoring for software systems relies on error-rate thresholds, latency
percentiles, and availability probes.  All three measure \emph{presence} or
\emph{absence} of response, not \emph{quality} or \emph{distributional
properties} of the response.  An LLM system can pass all three standard
monitoring criteria while simultaneously producing outputs that are
systematically wrong for a user cohort (C1b), returning stale retrieval results
(C3d), or operating with a broken evaluation pipeline that cannot detect either
(C4c).

Production LLM monitoring must extend to distributional monitoring: output
token-length distribution, semantic similarity to baseline outputs,
task-specific quality proxy metrics, and retrieval freshness timestamps.
The six failure classes in this taxonomy map to distinct monitoring primitives;
Table~\ref{tab:monitoring} provides a mapping.

\begin{table}[t]
\caption{Recommended monitoring primitives per failure class.}
\label{tab:monitoring}
\centering
\small
\resizebox{\linewidth}{!}{%
\begin{tabular}{lL{4.5cm}}
\toprule
\textbf{Class} & \textbf{Monitoring Primitive} \\
\midrule
C1 Drift         & Output distribution shift; semantic drift score; A/B to
                   baseline \\
C2 Infrastructure& Error rate; P99 latency; GPU memory utilisation; queue depth \\
C3 Integration   & System prompt integrity; context window utilisation;
                   retrieval freshness; tool-call success rate \\
C4 Evaluation    & Benchmark--production correlation; evaluation set
                   contamination check; per-cohort metric decomposition \\
C5 Safety        & PII regex on outputs; citation verification; policy
                   classifier; audit log completeness \\
C6 Operational   & Alert coverage matrix; runbook inventory; escalation
                   path test \\
\bottomrule
\end{tabular}}
\end{table}

\subsection{Safety \& Compliance Is Not a Model Problem}

C5 is the co-dominant failure class at 26\% of incidents, but its sub-classes
span the full spectrum from model behaviour (5a: PII memorisation) to system
design (5b: absent citation verification) to governance (5d: auditability gap).
The shared property is not model behaviour; it is the \emph{absence of a
verification layer} between model output and consequential action.  No amount
of alignment training eliminates the need for verification layers in FC\_A use
cases: even a model that hallucinates citations at a rate of 0.1\% --- well
within current state-of-the-art --- will produce a legal error every 1{,}000
requests without verification.  The taxonomy enables targeted remediation:
each C5 sub-class implies a distinct control.

\subsection{Evaluation Failure Is Underreported}

C4 accounts for only 10\% of incidents in our dataset --- the lowest share ---
but we believe this is a substantial undercount.  Evaluation failures are
meta-failures: they are the failure of the signal that would catch other
failures.  They rarely surface as reportable incidents.  No court document
records an LLM evaluation pipeline failing to detect benchmark contamination.
No regulatory filing identifies metric proxy collapse as the cause of customer
harm.  But both failure modes degrade the system's ability to detect and
respond to the other five failure classes.  Dedicated evaluation
infrastructure --- with production traffic sampling, regular benchmark rotation,
and per-cohort metric decomposition --- is a first-class engineering concern,
not an academic exercise.

\subsection{Failure Budget as Organisational Design Tool}

The failure budget framework is not only a monitoring concept; it is a forcing
function for cross-functional alignment.  Assigning a use case to FC\_A requires
legal, compliance, ML engineering, and product teams to agree on acceptable
risk before a single model is selected.  This conversation --- ``what failure
rate can we accept for this workflow, and what are the operational consequences
of a breach?'' --- rarely happens in practice.  Teams instead negotiate model
accuracy on a benchmark proxy, which answers a different question.  The formal
budget framework makes the right question legible and auditable.  Under DORA
\citep{dora2022} and the EU AI Act \citep{euaiact2024}, regulated entities are
required to document and test the resilience of critical information
infrastructure.  A failure budget framework provides the structured basis for
that documentation.

% =============================================================================
\section{Conclusion}
\label{sec:conclusion}
% =============================================================================

We presented the first system-level failure taxonomy for production LLM
deployments, comprising six failure classes that cover the full deployment
stack from model drift to operational process failure.  We validated the
taxonomy against 50 real-world incidents from verifiable public sources,
finding that 50\% of incidents are silent failures invisible to standard
monitoring, and that Safety \& Compliance and Integration failures are the
two dominant classes in publicly reportable incidents.

The failure budget framework provides a formal model for assigning acceptable
failure rates by use case risk class, shifting engineering teams from
benchmark-accuracy debates to system-level reliability reasoning.  The
companion reference implementation enables teams to instrument this framework
without building it from scratch.

Open problems include: inter-rater reliability study for the labeled dataset;
LLM-based classifier to complement the rule-based implementation; domain-specific
failure budget calibration for regulated sectors beyond financial services;
and temporal analysis of whether failure patterns are shifting as LLM deployment
matures and regulatory frameworks come into force.

All data, code, the taxonomy schema, the classifier implementation, and the
failure budget calculator are available at:
\url{https://github.com/priyanka25aug/llm-failure-taxonomy}.

% =============================================================================
% BIBLIOGRAPHY
% =============================================================================
\bibliographystyle{abbrvnat}
\bibliography{refs}

% =============================================================================
% APPENDICES
% =============================================================================
\appendix
\onecolumn

\section{Full Taxonomy Schema}
\label{appendix:schema}

The complete taxonomy definition, including all sub-classes, key indicators,
detectability profiles, and failure budget applicability, is maintained as a
machine-readable YAML file at
\url{https://github.com/priyanka25aug/llm-failure-taxonomy/blob/main/taxonomy/taxonomy.yaml}.
A JSON Schema for labeling incident records is provided at
\url{https://github.com/priyanka25aug/llm-failure-taxonomy/blob/main/taxonomy/schema.json}.

Table~\ref{tab:full_schema} reproduces the sub-class definitions.

\begin{table}[ht]
\caption{Complete taxonomy sub-class reference.}
\label{tab:full_schema}
\centering
\small
\begin{tabular}{llL{8.0cm}L{4.5cm}}
\toprule
\textbf{Class} & \textbf{Sub} & \textbf{Definition} & \textbf{Key Indicator} \\
\midrule
\multirow{3}{*}{C1 Drift}
  & 1a & Upstream provider silent model update & Output distribution shift \\
  & 1b & Production input distribution drift  & Semantic drift score \\
  & 1c & Context window saturation drift       & Avg.\ context utilisation $>$85\% \\
\midrule
\multirow{4}{*}{C2 Infra}
  & 2a & P99 latency regression                & Latency SLO breach \\
  & 2b & OOM / KV-cache exhaustion             & GPU memory $>$95\%; 5xx rate \\
  & 2c & Routing misclassification under load  & Wrong model serving traffic \\
  & 2d & API breaking change                   & Response schema mismatch \\
\midrule
\multirow{5}{*}{C3 Integration}
  & 3a & Prompt injection / jailbreak          & System prompt integrity check \\
  & 3b & Context truncation (silent)           & Context utilisation $>$90\% \\
  & 3c & Tool-call hallucination               & Tool-call success rate \\
  & 3d & RAG retrieval mismatch (stale index)  & Retrieval freshness timestamp \\
  & 3e & Multi-turn state corruption           & Per-session coherence score \\
\midrule
\multirow{4}{*}{C4 Evaluation}
  & 4a & Metric proxy collapse / benchmark gaming & Benchmark--production gap \\
  & 4b & Evaluation set contamination           & Training/eval overlap rate \\
  & 4c & Production--evaluation distribution gap & Covariate shift score \\
  & 4d & Point accuracy vs.\ distributional     & Per-cohort metric decomposition \\
\midrule
\multirow{5}{*}{C5 Safety}
  & 5a & PII leakage / memorisation             & PII regex on outputs \\
  & 5b & Hallucinated legal/medical citations   & Citation verification pipeline \\
  & 5c & Policy boundary violation              & Output policy classifier \\
  & 5d & Auditability gap                       & Audit log completeness \\
  & 5e & Copyright reproduction                 & Similarity to training corpus \\
\midrule
\multirow{4}{*}{C6 Operational}
  & 6a & Monitoring blind-spot                  & Alert coverage matrix \\
  & 6b & Runbook absence                        & Runbook inventory \\
  & 6c & Escalation breakdown                   & Time to escalation vs.\ SLO \\
  & 6d & Canary / shadow test gap               & Shadow traffic coverage \\
\bottomrule
\end{tabular}
\end{table}

\clearpage
\section{Dataset Sample}
\label{appendix:dataset}

Table~\ref{tab:dataset_sample} provides two representative incidents per
failure class.  Full 50-incident dataset with source URLs at
\url{https://github.com/priyanka25aug/llm-failure-taxonomy/blob/main/data/public_incidents/incidents.csv}.

\begin{table}[ht]
\caption{Representative incident sample (2 per class).}
\label{tab:dataset_sample}
\centering
\small
\resizebox{\linewidth}{!}{%
\begin{tabular}{llL{4.5cm}lll}
\toprule
\textbf{ID} & \textbf{Class} & \textbf{Incident} & \textbf{Sev.} &
\textbf{Detect.} & \textbf{MTTD} \\
\midrule
FTX-0001 & C1a & GPT-4 silent behaviour change (Mar 2023) & High & Silent & 168h \\
FTX-0002 & C1b & Production query distribution shift post-launch & Medium & Silent & 336h \\
\midrule
FTX-0007 & C2b & vLLM KV-cache exhaustion under bursty load & High & Immediate & 0.5h \\
FTX-0008 & C2a & ChatGPT outage latency spike (Dec 2022) & High & Immediate & 1h \\
\midrule
FTX-0013 & C3a & Bing Chat Sydney prompt injection (Feb 2023) & Critical & Delayed & 24h \\
FTX-0014 & C3d & RAG stale index serving outdated Basel III guidance & Critical & Silent & 2880h \\
\midrule
FTX-0026 & C4a & MMLU fine-tuned model degraded real-world performance & Medium & Silent & 720h \\
FTX-0027 & C4b & Evaluation set contamination in benchmark study & Medium & Silent & 1440h \\
\midrule
FTX-0031 & C5b & Mata v.\ Avianca hallucinated citations (2023) & Critical & Delayed & 48h \\
FTX-0032 & C5a & Samsung employee PII leak via ChatGPT (2023) & Critical & Delayed & 72h \\
\midrule
FTX-0044 & C6a & Output-length drift unmonitored for 3 weeks & Medium & Delayed & 504h \\
FTX-0045 & C6b & No runbook for model API deprecation scenario & High & Delayed & 120h \\
\bottomrule
\end{tabular}}
\end{table}

\clearpage
\section{Failure Budget Calculator}
\label{appendix:calculator}

The reference implementation computes budget utilisation per use case.
Core logic (Python pseudocode):

\begin{minipage}{\columnwidth}
\begin{Verbatim}[fontsize=\small]
BUDGET_RATES = {FC_A: 1.0, FC_B: 5.0,
                FC_C: 20.0, FC_D: 100.0}
SEVERITY_WEIGHTS = {critical: 3, high: 2,
                    medium: 1, low: 0.5}

def budget_utilisation(uc, failures):
  weighted = sum(SEVERITY_WEIGHTS[f.severity]
                 for f in failures)
  rate = weighted / (uc.requests / 1000)
  return (rate / BUDGET_RATES[uc.risk]) * 100
\end{Verbatim}
\end{minipage}

\vspace{0.5em}
Full implementation with reporting and multi-use-case portfolio tracking:
{\small\url{https://github.com/priyanka25aug/llm-failure-taxonomy/blob/main/src/failure_budget/calculator.py}}.

\end{document}
